\clearpage\section{METODOLOGIA}
Os procedimentos metodológicos usados em cada disciplina estão especificados, em linhas gerais, nos respectivos Programas de Unidade Didática (PUDs), mas dependem, adicionalmente, das características de cada professor. A grande maioria dos professores opta por aulas expositivas, conforme as necessidades de cada disciplina, com auxílio de quadro branco e pincel, intercaladas com o uso de projeções, aulas de exercícios, práticas em laboratórios, salas de informática, ou ainda visitas a setores do próprio campus ou externas a este. Recursos adicionais também estarão presentes, como o uso de ferramentas de simulação numérica (em determinadas disciplinas específicas) ou tecnologias que garantam a acessibilidade de docentes e/ou discentes (quando necessário).

Obedecendo a Resolução CNE/CES 11, de 11 de março DE 2002, as disciplinas do \NOMECURSO estão divididas em núcleo de conteúdos básicos, um núcleo de conteúdos profissionalizantes e um núcleo de conteúdos específicos. A descrição detalhada sobre cada núcleo é apresentada no item 10.1 que trata da Organização curricular.

Algumas disciplinas de conteúdo básico, como Introdução à Engenharia, HST, Ética e Projetos Sociais, assim como algumas disciplinas de gestão, respectivamente listadas nas Tabs. \ref{tab5} e \ref{tab6}, usarão diferentes metodologias, tais como: trabalho em equipe, seminários de apresentação de projetos por parte dos alunos, pesquisas diversas, trabalho de campo, convivência industrial no caso de estágios, entre outras atividades. Além disso, como já acontece nas outras disciplinas, o professor disponibiliza um horário extra-classe para tira-dúvidas e assim efetuar um melhor acompanhamento da aprendizagem dos alunos.

Tendo consciência que o perfil do egresso deve estar antenado com o mundo do trabalho em constante mutação, busca-se formar um profissional com o conhecimento específico de sua profissão, mas também com uma visão do todo. Este profissional deve saber buscar conhecimento a todo momento, ficar atento a novas tecnologias e desenvolvimentos, e possuir habilidades de comunicação com outras áreas, liderança, administração, entre outras. Neste contexto, o aluno através da interdisciplinaridade existente no curso tem a possibilidade de cursar várias disciplinas do \NOMECURSO em outras engenharias e/ou outras áreas distintas, favorecendo uma troca de experiências e uma visão mais ampla durante sua formação. Além disso, a interdisciplinaridade ocorre entre as mais variadas disciplinas do curso em que um conteúdo é útil para uma outra disciplina subsequente e assim sucessivamente. Neste sentido, foi inserido um campo nos PUDs e na matriz do curso para que seja destacada a interdisciplinaridade de cada disciplina, como forma de deixar claro ao aluno, bem como de incentivar o professor a buscar que esta seja efetivamente alcançada.

Os conteúdos pertinentes às políticas de educação ambiental (lei nº 9.795, de 27 de abril de 1999 e Decreto Nº 4.281 de 25 de junho de 2002), de educação em direitos humanos (Resolução Nº 1, de 30 de maio de 2012) e de educação das relações étnico-raciais (Resolução CNE/CP N° 1, de 17 de junho de 2004), conforme as normatizações vigentes, são contemplados em sua magnitude, da seguinte forma:
%Os conteúdos obrigatórios, recomendados pelas Resoluções CNE Nº2, de 15 de Junho de 2012, CNE Nº1, de 30 de maio de 2012 e CNE Nº1, de 17 de junho de 2004, são contemplados em sua magnitude da seguinte forma:

\begin{itemize}
	\item \textbf{Educação Ambiental:} Disciplina de Ética e Cidadania 
	\item \textbf{Relações Étnicas-Raciais:} Disciplina de Projetos Sociais 
	\item \textbf{Direitos Humanos:} Disciplina de Projetos Sociais
\end{itemize}

A forma de abordagem das temáticas mencionadas, além da abordagem enquanto conteúdos nos componentes curriculares citados, também poderá levar em consideração alguns aspectos, a saber: incentivo a pesquisas envolvendo as temáticas; desenvolvimento de projetos de extensão; organização de eventos, palestras e realização de visitas técnicas.

Nas disciplinas que envolvem práticas de laboratório, haverá um contato maior com os equipamentos didáticos, colocando o aluno em contato direto com os fenômenos físicos, envolvendo ainda recursos de informática para a aquisição e tratamento de dados, bem como para a confecção de relatórios. Em geral, o aluno de \NOMEENG dispõe de um grande arsenal de ferramentas de informática que vão auxiliar diretamente em seus estudos, juntamente com os recursos da Internet, da Biblioteca, bem como da Biblioteca Virtual (\href{http://bvu.ifce.edu.br/login.php}{http://bvu.ifce.edu.br/login.php}). 

Adicionalmente, um elevado número de alunos faz parte dos corpos técnicos dos laboratórios e grupos de pesquisa, dispondo de computadores, envolvimento com trabalhos correlatos, que permitem um melhor acompanhamento das disciplinas cursadas naquele momento. 

Finalmente, o uso das TIC's (Tecnologia da Informação e Comunicação) oferecem um conjunto de recursos tecnológicos no processo de ensino-aprendizagem, tais como e-mail, grupos de discussão on-line e redes sociais em salas de vídeo conferência que possibilitam o intercâmbio de informações e a geração de novos conhecimentos e competências entre todos os envolvidos. Além disso, estas também garantem a acessibilidade de docentes e/ou discintes com necessidades especiais.

Atividades de monitoria são realizadas sob orientação de um professor orientador para alunos que estejam com dificuldade de aprendizagem e, assim, contribuir para um maior envolvimento dos alunos com o IFCE, propiciando uma melhor formação acadêmica ao aluno, além de estimulá-los à participação no processo educacional e nas atividades relativas ao ensino.

Durante o processo de ensino-aprendizagem e seguindo as orientações legais dispostas nos Decretos N° 5.296/2004, N° 6.949/2009, N° 7.611/2011, na Portaria N° 3.284/2003, na Lei N° 12.764, de 27 de dezembro de 2012, são identificados possíveis estudantes com necessidades especiais, os quais, terão um tratamento diferenciado, com o devido apoio ao discente prestado por uma equipe multidisciplinar do campus.   

Outras atividades de apoio ao discente são comentadas no tópico \ref{APOIO}, APOIO AO DISCINTE.
